\chapter*{Введение}
\addcontentsline{toc}{chapter}{Введение}

Практическая работа посвящена статическому и динамическому анализу приложения
поиска палиндромных подстрок. Приложение реализовано на JavaScript в среде
Node.js и использует алгоритм Манакера как основной алгоритм решения наукоёмкой
задачи обработки строк.

Цель работы --- выполнить фаззинг-тестирование и статический анализ приложения,
зафиксировать найденные дефекты, сопоставить предупреждения анализаторов с CWE и
БДУ ФСТЭК России, а также оценить удобство и эффективность использованных
инструментов.

Для достижения цели решены следующие задачи:
\begin{enumerate}
    \item выбраны два модуля приложения, в которые внесены две ошибки для анализа;
    \item выполнено фаззинг-тестирование этих модулей с использованием Jazzer.js;
    \item описаны команды запуска фаззера, настройки окружения и смысл параметров;
    \item зафиксированы входные данные, позволившие выявить ошибки;
    \item выбраны два статических анализатора: SonarJS-совместимые правила ESLint и ESLint Security;
    \item выполнен статический анализ кода приложения;
    \item предупреждения анализаторов сопоставлены с CWE и БДУ ФСТЭК России;
    \item составлены итоговые таблицы и выводы по результатам динамического и статического анализа.
\end{enumerate}
